\kchapter{ConfigMaps and Secrets}{cm}

\begin{objectives}
\begin{itemize}
  \item why configuration must live outside the image;
  \item creating ConfigMaps and consuming them as env vars or files;
  \item Secrets: types, consumption, and the truth about base64;
  \item how configuration changes reach running Pods;
  \item production secret management: encryption at rest, External Secrets,
        Sealed Secrets, SOPS.
\end{itemize}
\end{objectives}

\section{Config belongs outside the image}

The same image should run in development, staging, and production. Only its
configuration should differ: database host, log level, feature flags, API keys.
This is principle III of the well-known ``Twelve-Factor App''. Kubernetes gives
you two objects for it:

\begin{itemize}
  \item \kicon{cm}~\kterm{ConfigMap} --- non-sensitive settings;
  \item \kicon{secret}~\kterm{Secret} --- passwords, tokens, keys, certificates.
\end{itemize}

\section{ConfigMaps}

\begin{yamlcode}{app-config.yaml}
apiVersion: v1
kind: ConfigMap
metadata:
  name: app-config
data:
  LOG_LEVEL: "info"
  FEATURE_SEARCH: "true"
  DB_HOST: "postgres"
  app.properties: |              # a whole file as one key
    cache.size=256
    cache.ttl=300s
\end{yamlcode}

Or create one from the command line:

\begin{bashcode}{terminal}
kubectl create configmap app-config \
  --from-literal=LOG_LEVEL=info \
  --from-file=app.properties \
  --dry-run=client -o yaml > app-config.yaml
\end{bashcode}

\subsection{Consuming a ConfigMap}

\begin{yamlcode}{consume-config.yaml}
spec:
  containers:
    - name: api
      image: ghcr.io/example/api:1.4.0
      # 1) one key -> one env var
      env:
        - name: LOG_LEVEL
          valueFrom:
            configMapKeyRef: { name: app-config, key: LOG_LEVEL }
      # 2) every key -> env vars
      envFrom:
        - configMapRef: { name: app-config }
      # 3) keys -> files in a directory
      volumeMounts:
        - name: config
          mountPath: /etc/app        # /etc/app/app.properties, /etc/app/LOG_LEVEL ...
          readOnly: true
  volumes:
    - name: config
      configMap:
        name: app-config
        items:                       # optional: pick specific keys
          - key: app.properties
            path: app.properties
\end{yamlcode}

\subsection{How changes propagate}

\begin{center}
\ktablestart
\begin{tabularx}{\linewidth}{P{4.0cm} L}
\kheader \kth{Consumed as\ldots} & \kth{After you edit the ConfigMap\ldots} \\
Environment variables & \textbf{Nothing changes} until the Pod restarts. Run \kc{kubectl rollout restart deployment/api}. \\
Mounted volume & Files update automatically within about a minute (not with \kc{subPath} mounts). The app must re-read them. \\
\end{tabularx}
\end{center}

\begin{protip}
A popular pattern is to put a hash of the config into a Pod-template
annotation, so that any config change produces a new template and therefore an
automatic rolling update. Kustomize's \kc{configMapGenerator} and Helm's
\kc{sha256sum} function do exactly this (Chapter 18).
\end{protip}

Mark ConfigMaps or Secrets \kc{immutable: true} when they should never change;
this protects against accidents and reduces API-server load.

\section{Secrets}

\begin{yamlcode}{db-secret.yaml}
apiVersion: v1
kind: Secret
metadata:
  name: db-secret
type: Opaque
stringData:                     # plain text here; stored base64-encoded
  username: app
  password: "s3cr3t-Pa55w0rd"
\end{yamlcode}

\begin{bashcode}{terminal}
kubectl create secret generic db-secret \
  --from-literal=username=app --from-literal=password='s3cr3t-Pa55w0rd'

kubectl get secret db-secret -o jsonpath='{.data.password}' | base64 -d
\end{bashcode}

\begin{warningbox}
\textbf{base64 is encoding, not encryption.} Anyone who can read the Secret
object can read the value. Never commit plain Secret manifests to Git. Protect
Secrets with RBAC, enable encryption at rest, and use one of the tools in
Section~\ref{sec:secret-tools}.
\end{warningbox}

\subsection{Secret types}

\begin{center}
\ktablestart
\begin{tabularx}{\linewidth}{P{5.6cm} L}
\kheader \kth{Type} & \kth{Use} \\
\kc{Opaque} & Arbitrary key/value data (default). \\
\kc{kubernetes.io/tls} & A certificate and key (\kc{tls.crt}, \kc{tls.key}) for Ingress/Gateway. \\
\kc{kubernetes.io/dockerconfigjson} & Credentials for pulling from private registries. \\
\kc{kubernetes.io/basic-auth}, \kc{ssh-auth} & Structured credentials. \\
\end{tabularx}
\end{center}

\begin{bashcode}{terminal}
# A registry pull secret, then reference it from a Pod or ServiceAccount
kubectl create secret docker-registry ghcr-pull \
  --docker-server=ghcr.io --docker-username=USER --docker-password=TOKEN

kubectl create secret tls shop-tls --cert=tls.crt --key=tls.key
\end{bashcode}

\begin{yamlcode}{use-secret.yaml}
spec:
  imagePullSecrets:
    - name: ghcr-pull
  containers:
    - name: api
      image: ghcr.io/example/private-api:1.4.0
      env:
        - name: DB_PASSWORD
          valueFrom:
            secretKeyRef: { name: db-secret, key: password }
      volumeMounts:
        - name: creds
          mountPath: /var/run/secrets/app
          readOnly: true
  volumes:
    - name: creds
      secret:
        secretName: db-secret
        defaultMode: 0400
\end{yamlcode}

\begin{protip}
Prefer mounting Secrets as \emph{files} over env vars. Environment variables
leak easily (crash dumps, \kc{/proc}, child processes, debug endpoints that print
the environment).
\end{protip}

\section{Production secret management}
\label{sec:secret-tools}

\begin{center}
\ktablestart
\begin{tabularx}{\linewidth}{P{3.6cm} L}
\kheader \kth{Approach} & \kth{How it works} \\
Encryption at rest & Configure the API server (or your cloud's KMS integration) so Secrets are encrypted inside etcd. Managed clusters usually offer a checkbox. \\
External Secrets Operator & Secrets live in AWS Secrets Manager, GCP Secret Manager, Azure Key Vault, HashiCorp Vault\ldots An \kc{ExternalSecret} object tells the operator to sync them into Kubernetes Secrets. \\
Sealed Secrets & Encrypt a Secret with the cluster's public key into a \kc{SealedSecret} that is safe to commit. Only the in-cluster controller can decrypt it. \\
SOPS (+ age/KMS) & Encrypt values in YAML files in Git; decrypt in your CD pipeline or with a Kustomize/Helm plugin. \\
Secrets Store CSI Driver & Mount secrets from an external vault directly as files, without creating Kubernetes Secrets. \\
\end{tabularx}
\end{center}

\begin{yamlcode}{external-secret.yaml}
apiVersion: external-secrets.io/v1
kind: ExternalSecret
metadata:
  name: db-secret
spec:
  refreshInterval: 1h
  secretStoreRef:
    name: aws-secrets-manager
    kind: ClusterSecretStore
  target:
    name: db-secret               # the Kubernetes Secret to create
  data:
    - secretKey: password
      remoteRef:
        key: prod/taskflow/db
        property: password
\end{yamlcode}

\begin{recap}
\begin{itemize}
  \item ConfigMaps for settings, Secrets for credentials --- never bake either into images.
  \item Consume as env vars (need restart) or files (update live).
  \item Secrets are only base64-encoded; protect them with RBAC, encryption at
        rest, and an external secret manager.
\end{itemize}
\end{recap}
